% Problem Set 1 -- 17 problems. Do not edit except to fix statement typos.
\problemset{1}{}

\problem{1}{1}{Let $A$ be a set and $\Pow(A)$ be its power set. Prove that there is no
injective map $f\colon \Pow(A)\to A$ and no surjective map $g\colon A\to\Pow(A)$.}

\problem{1}{2}{Let $X$ be an infinite set. Is the collection
$\T_\infty := \{U\subseteq X \mid X\setminus U \text{ is infinite or } U=\emptyset\}$
a topology on $X$? Justify your answer.}

\problem{1}{3}{Let $X=\{a,b,c\}$. Determine which of the following collections of subsets
form a valid topology on $X$. For those that fail, explicitly state which axiom or axioms
are violated:
\begin{enumerate}[label=(\roman*)]
  \item $\T_1=\{\emptyset,\{a\},\{b\},X\}$;
  \item $\T_2=\{\emptyset,\{a\},\{a,b\},X\}$;
  \item $\T_3=\{\emptyset,\{a,b\},\{b,c\},X\}$.
\end{enumerate}}

\problem{1}{4}{Let $X$ be a topological space and $A\subset X$. Prove that $A$ is open in
$X$ if and only if for each $x\in A$, $\exists$ an open set $U$ in $X$ such that
$x\in U\subseteq A$.}

\problem{1}{5}{Let $U$ be an open set and $V$ be a closed set in $X$. Prove that
$U\setminus V$ is open in $X$ and $V\setminus U$ is closed in $X$.}

\problem{1}{6}{Let $X=\N$ and define $\T_1,\T_2\subseteq\Pow(\N)$ as follows:
\begin{enumerate}[label=(\roman*)]
  \item $\T_1 := \{\emptyset,\N,\{1\},\{1,2\},\{1,2,3\},\{1,2,3,4\},\dots\}$.
  \item $\T_2 := \{\emptyset,\N,\N\setminus\{1\},\N\setminus\{1,2\},\N\setminus\{1,2,3\},
        \N\setminus\{1,2,3,4\},\dots\}$.
\end{enumerate}
Determine whether $\T_1$ and $\T_2$ define topologies on $\N$.}

\problem{1}{7}{Show that if $\B$ is a basis for a topology on $X$, then the topology
generated by $\B$ equals the intersection of all topologies on $X$ that contain $\B$.}

\problem{1}{8}{Let $X$ be an infinite set and $\T_f$ be the finite complement topology on
$X$. If $U$ and $V$ are non-empty open sets in $(X,\T_f)$, prove that $U\cap V\neq\emptyset$.}

\problem{1}{9}{Let $X\neq\emptyset$ be a set and let $a\in X$ be a fixed point. Prove that
the \emph{particular point} collection
\[ \T_a=\{V\subseteq X \mid V=\emptyset \text{ or } a\in V\} \]
defines a valid topology on $X$.}

\problem{1}{10}{Let $X\neq\emptyset$ be a set and let $a\in X$ be a fixed point. Prove that
the \emph{excluded point} collection
\[ \T^a := \{V\subseteq X \mid V=X \text{ or } a\notin V\} \]
defines a valid topology on $X$.}

\problem{1}{11}{Let $X$ be an infinite set and let $a\in X$ be a fixed point. Define a
collection of subsets $\T_a\subseteq\Pow(X)$ by:
\[ \T_a=\Pow\bigl(X\setminus\{a\}\bigr)\cup\{V\subseteq X \mid a\in V \text{ and }
   X\setminus V \text{ is finite}\}. \]
Prove that $\T_a$ defines a valid topology on $X$. What happens when $X$ is a finite set?}

\problem{1}{12}{Let $X$ be an infinite set and let $a,b\in X$ be two distinct points.
Define a collection of subsets $\T\subseteq\Pow(X)$ by:
\[ \T=\Pow(X\setminus\{a,b\})\cup\{V\subseteq X \mid a\in V \ \text{ or }\ b\in V
   \ \text{ and }\ X\setminus V \text{ is finite}\}. \]
Is $\T$ a topology on $X$?}

\problem{1}{13}{Let $\B_1$, $\B_2$, and $\B_3$ be respectively the collections of all open
balls, open rectangles, and open triangles in the plane $\R^2$.
\begin{enumerate}[label=(\roman*).]
  \item Prove that all three collections satisfy the basis axioms on $\R^2$.
  \item Let $\T_i$ be the topology generated by the basis $\B_i$ for $i=1,2,3$. Let
        $f\colon(\R^2,\T_i)\to(\R^2,\T_j)$ be the identity map. Prove that for any choice
        of $i,j\in\{1,2,3\}$, the map $f$ is continuous.
  \item Conclude that $\T_1=\T_2=\T_3$.
\end{enumerate}}

\problem{1}{14}{Along with the finite complement topology ($\T_f$) and the countable
complement topology ($\T_c$), we consider the topologies generated by the following six bases:
\begin{enumerate}[label=(\roman*)]
  \item \textbf{Standard Topology ($\T_{\mathrm{std}}$):} Generated by the basis of all open
        intervals: \[\B_1=\{(a,b)\mid a<b \text{ and } a,b\in\R\}.\]
  \item \textbf{$K$-Topology ($\T_K$):} Generated by the basis:
        \[\B_2=\{(a,b)\mid a<b\}\cup\{(a,b)\setminus K\mid a<b\},\quad\text{where }
          K=\{1,\tfrac12,\tfrac13,\dots\}.\]
  \item \textbf{Lower Limit Topology ($\T_\ell$):} Generated by the basis of all half-open
        intervals closed on the left: \[\B_3=\{[a,b)\mid a<b \text{ and } a,b\in\R\}.\]
  \item \textbf{Upper Limit Topology ($\T_u$):} Generated by the basis of all half-open
        intervals closed on the right: \[\B_4=\{(a,b]\mid a<b \text{ and } a,b\in\R\}.\]
  \item \textbf{Left-Ray Topology ($\T_{\mathrm{left}}$):} Generated by the basis of open
        left rays: \[\B_5=\{(-\infty,b)\mid b\in\R\}.\]
  \item \textbf{Right-Ray Topology ($\T_{\mathrm{right}}$):} Generated by the basis of open
        right rays: \[\B_6=\{(a,\infty)\mid a\in\R\}.\]
\end{enumerate}
For each pair of these topologies, determine whether one contains the other (i.e., whether
one is finer, coarser, or if they are completely incomparable). Summarize your results by
drawing a partial order hierarchy diagram representing the containment relations among all
eight topologies.}

\problem{1}{15}{Let $\mathcal{T}_X$ denote the collection of all topologies on a non-empty
set $X$ and let $\mathcal{T}_X/\mathrm{hom}$ be the set of homeomorphism classes of topologies
on $X$. Let $S_X$ be the group of all bijective maps on $X$. Note that each bijection
$f\colon X\to X$ induces a bijection $\tilde f\colon\Pow(X)\to\Pow(X)$ defined by
$\tilde f(A)=\{f(a)\mid a\in A\}$ for every subset $A\in\Pow(X)$. Consequently, $S_X$ acts
naturally on $\mathcal{T}_X$ via $f.\T=\tilde f(\T)$ for all $f\in S_X$ and
$\T\in\mathcal{T}_X$. Prove that there is a natural bijection between
$\mathcal{T}_X/\mathrm{hom}$ and the collection of orbits of the natural action of $S_X$ on
$\mathcal{T}_X$.}

\problem{1}{16}{List down all possible (there are 29) topologies on a three-point set
$X=\{a,b,c\}$. Compute the number of orbits under the natural action of $S_3$ on the set
$\mathcal{T}_X$ of all topologies on $X$.}

\problem{1}{17}{Let $X$ be a set, and let $\mathcal{S}_1$ and $\mathcal{S}_2$ be two
subbases on $X$ generating topologies $\T_1$ and $\T_2$, respectively.
\begin{enumerate}[label=(\roman*).]
  \item If $\mathcal{S}_1\subseteq\mathcal{S}_2$, prove that $\T_1\subseteq\T_2$.
  \item Give an explicit example where $\T_1=\T_2$ even though
        $\mathcal{S}_1\neq\mathcal{S}_2$.
\end{enumerate}}
